\documentclass[10pt]{article}

\def\Type{LessonDJ}
\def\TypeNum{Workshop Week 9} %Lesson #
\def\TypeTitle{Optimization and Net Change } %Lesson Title
\def\Semester{Fall 2021} %Not Used 
\def\Course{MATH 1190}%Course number and name
\def\Targets{  \target{A3}, \target{I1} }

\usepackage{preambleDJ} 

%%%%%%%%%%%% BEGIN DOCUMENT %%%%%%%%%%%%%%%
%%%%%%%%%%%% BEGIN DOCUMENT %%%%%%%%%%%%%%%
%%%%%%%%%%%% BEGIN DOCUMENT %%%%%%%%%%%%%%%
%%%%%%%%%%%% BEGIN DOCUMENT %%%%%%%%%%%%%%%
%%%%%%%%%%%% BEGIN DOCUMENT %%%%%%%%%%%%%%%

\begin{document}
\pagestyle{otherpages}
\thispagestyle{firstpage}


\vs{.6}
\
%\begin{center}
%    \includegraphics[width=0.8\textwidth]{IMG_2891.jpeg}
%\end{center}




\section{Optimization of the product of two numbers}
\begin{enumerate}
\item Find two numbers with a sum of $50$ that maximizes the product of the two numbers. Let's use , $x$ and $y$ to represent our two numbers.
%\[
%f(x,y) = xy
%\]
%subject to the constraint
%\[
%x^2 + 4y^2 = 16.
%\]

{\bf Modeling Phase}
\tasksA{2}{
\task Write down a formula for the product, $P$ of our two numbers ($x$ and $y$)
\task Write down a formula that represents ``The sum of two numbers is $50$." This should involve both $x$ and $y$.  This formula is called a {\em constraint}.
\task Use the constraint to express one variable in terms of the other.
\task Substitute your expression into the product, $P$ to obtain a single-variable function.
\task Simplify \( P \) and determine the domain of possible \(x\)-values.
}

\vs{1}

{\bf Calculus Phase}
\\ Now that you have obtained $P$ as a single-variable function, use the usual techniques to find the maximum and the minimum of $P$ (if any). 


\newpage
\section{Net Change and The Definite Integral}
\item Consider the graph of a function $f(t)$  which represents the velocity (in m/s) of a braking vehicle (looks like one of those ``rolling stops") over time (in seconds s). 
%In addition, suppose this function is given by $f(x) = 10\log\left(2\pi\ x\right).$
\graphicsC{bus_go_broooom.png}{4}

%Consider the following questions:
    \begin{enumerate}
    \item What does the definite integral $\ds \int_5^{15} f(t) dt$ represent? In particular, what do the bounds symbolize \textbf{ and } how would you explain what the integral calculates (give units).\\
        \vs{1.5}
    \item Usie two different approaches to estimate the  area between the graph and the horizontal axis from $t = 5$ to $t=15$. What does this quantity represent (use units)?
    \enum{
    \item Use at least $2$ rectangles and a left endpoint sum.
    \item Use at least $2$ rectangles and a right endpoint sum.
    }
\vfill

\item If the car is $40$m from the start at $t=5$ seconds, how far from the start is the car at $t=15$ seconds?
\vfill
    \end{enumerate}

\newpage
\section{Optimizing a run/swim}
\item Suppose Jack is standing next to a river that is 1 mile wide. His house is on the opposite of the river 3 miles down. Jack can swim at 5 miles per hour but he can run at 8 miles per hour. What is the shortest amount of time it can take Jack to reach his house? You will find this by finding the best choice for $d$ where Jack stops running and starts swimming.\\
\mpageT{.4}{.6}[\hs{.1}]{
{\bf Modeling Phase} We first need to develop a model of the total time: $$T_{total}=T_{run}+T_{swim}$$
The total horizontal distance to the house is $3$ miles.  Once Jack starts to swim, the remaining horizontal distance is $d$ miles (see Figure 1). 
}{\graphicsC{runSwim}{4}{Figure 1}
}
\tasksA{2}{
\task What is the distance Jack ran? Your answer should contain $d$.
\task How long does it take for Jack to run this distance? This is $T_{run}$.  Your answer should contain $d$.\vs{.5}
\task What is the distance Jack swam?  {\em Hint}: You will want to utilize Pythagorean's Theorem to find this distance. Your answer should contain $d$.
\task How long does it take for Jack to swim this distance? This is $T_{swim}$.  Your answer should contain $d$. \vs{.75}
\task* What is $T_{total}$? Your answer will contain $d$\vs{.25}
}
{\bf Calculus Phase} Think about the process you have been using to optimize quantities. Use that process to find the minimum time it takes for Jack to get home. You will definitely want to use technology for the last part of finding critical numbers.


    \vfill



\newpage

\section{Optimizing an athletic field}

\mpageT{.5}{.6}{
\item An athletic field is shaped like a rectangle with a semicircular track on each of its shorter sides (see Figure 1).  
The total perimeter of the track is fixed at $400$ meters.   The variable $R$ represents the radius  of each semi-circle.  The variable $L$ is the length  of the straight portion of the track.\\\\
Find the dimensions that enclose the largest possible area.
}{ \centering
\begin{tikzpicture}[scale=1.2, thick]
  % Parametri
  \def\R{1.5} % raggio dei semicerchi
  \def\L{4}   % lunghezza del rettangolo
%
  % Parte rettangolare (linee orizzontali)
  \draw (0,\R) -- (\L,\R);
  \draw (0,-\R) -- (\L,-\R);
  \draw[dotted] (0,\R) -- (0,0);
  \draw[dotted] (\L,\R) -- (\L,0);
  \node[circle, fill=black, inner sep=1pt] at (0,0) {};
  \node[circle, fill=black, inner sep=1pt] at (\L,0) {};
%
%
  % Semicerchi alle estremità
  \draw (0,-\R) arc[start angle=270, end angle=90, radius=\R];
  \draw (\L,\R) arc[start angle=90, end angle=-90, radius=\R];
%
  % Etichette
  \node at (\L/2, \R-0.4) {$L$};
  \node at (-\R+1.3, \R/2) {$R$};
  \node at (\L+\R-1.3, \R/2) {$R$};
  \node at (\L/2, -\R-0.4) {};
  %\node at (-0.3, \R+0.8) {Figure 1: Track shape};
%
\end{tikzpicture}
\vs{-.3}\captionof{figure}{Track Shape}
%\captionof{figure}{Track Shape}
}

{\bf Modeling Phase}
\tasksA{2}{
\task Write an expression for the total area enclosed by the track in terms of $L$ and $R$. Recall that the area of a half-circle is $A=\pi R^2/2$
\task Write an expression for the total perimeter in terms of your variables, and set it equal to $400$. Recall that the perimeter of a half-circle in $P=\pi R$\vs{.5}
\task Use the perimeter constraint to express one variable in terms of the other, and substitute it into the area formula.
}



\vs{.5}

{\bf Calculus Phase}
 Think about the process you have been using to optimize other quantities. Use that process here to find the dimensions that maximize the area of the athletic field.












\end{enumerate}

\end{document}